\unit{Exposé 236 -- Les schémas de Picard : propriétés générales}
\sid{236}
\src{221}
\seminar{14\textsuperscript{e} année, 1961/62, n\textsuperscript{o} 236\\Mai 1962}
\paperhead{TECHNIQUE DE DESCENTE ET THÉORÈMES D'EXISTENCE EN GÉOMÉTRIE ALGÉBRIQUE\\
VI. LES SCHÉMAS DE PICARD : PROPRIÉTÉS GÉNÉRALES}%
  {par Alexander GROTHENDIECK}

\fsection{0}{Compléments à l'exposé précédent ([3], V)\textsuperscript{(1)}}

Il y a eu quelques progrès concernant les questions d'existence de préschémas
de Picard soulevées dans [3], V :

\medskip
\noindent a. (Mumford). Il n'est pas vrai en général que lorsque
$f:X\longrightarrow S$ est un morphisme projectif et séparable (= plat à
fibres séparables), le préschéma de Picard $\operatorname{Pic}_{X/S}$ existe,
même si les fibres de $f$ sont de dimension 1, et si $S$ est le spectre d'un
anneau de valuation discrète complet. Un exemple en est fourni en prenant
$S=\operatorname{Spec}\mathbf R[[t]]$, et pour $X$ le sous-schéma de
$\mathbf P^2_S$ défini par l'équation (où $x,y,z$ sont des variables
homogènes)
\[
x^2+y^2=tz^2,
\]
qui représente donc une conique dégénérant en deux droites concourantes
géométriquement, mais la fibre spéciale sur le corps $\mathbf R$ étant
néanmoins irréductible (elle est donnée par l'équation $x^2+y^2=0$ sur
$\mathbf R$). On voit facilement que, après extension étale
$S'\longrightarrow S$, avec $S'=\operatorname{Spec}\mathbf C[[t]]$, le
préschéma de Picard de $X'/S'$ existe, et on en obtient une description
explicite comme somme de copies $\widetilde S'$, où $\widetilde S'$ est déduit de
$S'$ en dédoublant une infinité de fois l'origine. On constate aisément que la
donnée de descente sur $\operatorname{Pic}_{X'/S'}$ pour
$S'\longrightarrow S$ (donnée ici par les opérations du groupe de Galois
$\mathbf Z/2\mathbf Z$ de $S'$ sur $S$) n'est pas effective, le groupe
échangeant entre eux certains points dédoublés (de sorte qu'il y a des orbites
qui ne sont pas contenues dans un ouvert affine). Cependant, Mumford peut
montrer que si $f:X\longrightarrow S$ est un morphisme projectif séparable tel
que, pour tout $s\in S$, les composantes irréductibles de $X_s$ sont
géométriquement irréductibles relativement à $k(s)$, alors
$\operatorname{Pic}_{X/S}$ existe ; sa démonstration s'appuie sur un
raffinement de son théorème de passage au quotient, cf. [7]. Noter d'autre part
qu'il reste possible que, sans hypothèses sur les composantes irréductibles des
fibres $X_s$, le schéma $\operatorname{Pic}^{\tau}_{X/S}$ (qui sera introduit
plus bas) existe pourtant.

\src{222}
\noindent b. (Murre). Soit $X$ un schéma propre sur un corps, alors
$\operatorname{Pic}_{X/k}$ existe. La démonstration reprend en partie la
démonstration de Chevalley [1], et s'appuie à fond sur la structure de groupe
du foncteur de Picard.

Pour certains compléments concernant la théorie des schémas de Picard,
notamment en relation avec les schémas abéliens, il y aura intérêt à consulter
[7]. Enfin, une lacune notable du présent exposé est l'absence de « critères
d'équivalence », permettant de comparer le schéma de Picard d'un schéma
projectif, et de ses sections hyperplanes ; les théorèmes-clefs pour développer
de tels critères se trouvent dans [5], auquel il faut joindre les théorèmes
d'existence des schémas de Picard.

\fsection{1}{Propriétés topologiques des préschémas en groupes commutatifs}

Soient d'abord $k$ un corps, $G$ un préschéma en groupes sur $k$. Comme
l'élément neutre $e$, étant rationnel sur $k$, est nécessairement fermé, il
s'ensuit aussitôt que la diagonale de $G\times_kG$ est fermée, donc $G$ est
séparé : tout préschéma en groupes sur un corps est séparé. Nous désignerons
par $G^\circ$ la composante connexe de l'élément neutre $e$. Comme $e$ est
rationnel sur $k$, $G^\circ$ est en fait géométriquement connexe, i.e. la
formation de $G^\circ$ est compatible avec le changement de corps de base. Il
s'ensuit aussi que $G^\circ$ est stable par multiplication (ensemblistement),
et lorsque $G$ est localement noethérien, donc $G^\circ$ ouvert, on peut
considérer $G^\circ$ comme un sous-groupe ouvert de $G$. Dans la suite, nous
supposerons $G$ localement de type fini sur $k$ ; alors $G^\circ$ est
géométriquement irréductible et de type fini sur $k$. En effet, on peut
supposer $k$ algébriquement clos, donc de plus $G$ réduit (car
$G_{\mathrm{réd}}$ sera alors un sous-groupe de $G$, compte tenu que
$G_{\mathrm{réd}}\times_kG_{\mathrm{réd}}$ sera encore réduit), donc simple
sur $k$ sur un ouvert non vide, donc partout comme on voit en translatant ledit
ouvert. Mais alors $G$ est localement irréductible, donc ses composantes
irréductibles sont identiques à ses composantes connexes, donc $G^\circ$ est
irréductible. Soit alors $U$ un voisinage affine de $e$ dans $G^\circ$ ;
utilisant le fait que $G^\circ$ est irréductible, on voit tout de suite que
$U\cdot U=G^\circ$, ce qui prouve que $G^\circ$ est quasi-compact, donc de type
fini sur $k$.

Supposons pour simplifier que $G$ soit commutatif. Pour tout entier $n>0$, soit
$G^{(n)}$ l'image inverse de $G^\circ$ par l'homomorphisme $\varphi_n$ de
puissance $n$-ième dans $G$, donc $G^{(n)}$ est un sous-groupe ouvert de $G$.
Nous poserons :
\[
G^\tau=\bigcup_nG^{(n)}.
\]

\src{223}
\[
G^\sigma=\bigcup_{(n,p)=1}G^{(n)},
\qquad
G^\rho=\bigcup_hG^{(p^h)},
\]
où $p$ est l'exposant caractéristique pour le corps $k$. On obtient donc autant
de sous-groupes ouverts de $G$, satisfaisant
\[
G^\sigma\cap G^\rho=G^\circ,
\qquad
G^\sigma G^\rho=G^\tau.
\]

\result{REMARQUE}{} On peut construire le schéma en groupes quotient
$G/G^\circ=N$ (cf. [3], IV), et définir alors $G^\tau$, $G^\sigma$, $G^\rho$
comme les images inverses dans $G$ du sous-groupe de torsion de $N$ (resp. du complémentaire naturel de sa
composante $p$-primaire, somme des composantes $q$-primaires pour $q$
nombre premier $\neq p$, resp. de sa composante $p$-primaire). On
notera à ce propos que $N$ est un schéma en groupes discret séparable sur $k$,
donc (une fois choisi une clôture algébrique $\overline k$ de $k$, donnant lieu
à un groupe de Galois $\pi$) s'identifie à un groupe discret ordinaire sur
lequel $\pi$ opère par automorphismes. C'est de cette façon qu'on peut
interpréter de façon évidente la construction du sous-groupe de torsion et la
décomposition de ce dernier en ses composantes $q$-primaires. Lorsque $G$ est
le schéma de Picard d'un schéma $X$ propre sur $k$, $N$ pourra s'appeler le
\emph{schéma de Néron--Severi} (réduit) de $X$ sur $k$. Lorsque
$G^\circ_{\mathrm{réd}}$ est un sous-schéma en groupes de $G^\circ$, ce qui a
lieu en particulier chaque fois que $k$ est parfait ou que $G^\circ$ est
propre sur $k$ (par exemple $X$ géométriquement normal), il y a lieu également
d'introduire le quotient $N'=G/G^\circ_{\mathrm{réd}}$, qui a tendance à se
comporter mieux que $N$ du point de vue spécialisation, i.e. quand $X$ varie
dans une famille de schémas algébriques.

Soient maintenant $S$ un préschéma localement noethérien, et $G$ un préschéma
en groupes sur $S$, localement de type fini sur $S$. Nous ne supposons pas $G$
de type fini sur $S$, ni séparé sur $S$. Nous poserons alors
\[
G^\circ=\bigcup_{s\in S}(G_s)^\circ,
\]
et lorsque $G$ est commutatif :
\[
G^\tau=\bigcup_{s\in S}(G_s)^\tau,
\qquad
G^\sigma=\bigcup_{s\in S}(G_s)^\sigma,
\qquad
G^\rho=\bigcup_{s\in S}(G_s)^\rho.
\]

\src{224}
Ce sont là des parties de $G$, stables par la multiplication de $G$, ce qui
n'implique évidemment pas qu'elles puissent être définies à l'aide de
sous-préschémas en groupes de $G$. Notamment, il semble qu'il n'existe pas en
général de sous-préschéma en groupes de $G$ dont l'ensemble sous-jacent soit
$G^\circ$. Bien entendu, si l'un de ces ensembles est ouvert, alors, muni de la
structure induite, c'est un sous-préschéma en groupes ouvert de $G$. Nous
verrons qu'il en est toujours ainsi de $G^\tau$ ; ainsi, du point de vue des
foncteurs représentables, en particulier du point de vue « spécialisations »,
l'équivalence numérique se comporte de façon plus satisfaisante que
l'équivalence algébrique. Voici les principales propriétés générales des
ensembles qu'on vient de définir :

\result{THÉORÈME}{1.1} $G^\circ$, $G^\tau$, $G^\sigma$, $G^\rho$ sont
localement constructibles. De plus :

\noindent (i) $G^\circ$ est quasi-compact sur $S$. Lorsque les $G_s^\circ$
sont propres et $G$ séparé sur $S$, alors $G^\circ$ est propre sur $S$, donc
fermé dans $G$.

\noindent (ii) $G^\tau$ est ouvert. Si $G^\circ$ est fermé, il en est de même
de $G^\tau$.

\noindent (iii) Si $G^\circ$ est fermé, il en est de même de $G^\sigma$ pourvu
qu'on soit en égale caractéristique, i.e. tous les corps résiduels de $S$ ont
même caractéristique. Si $G^\circ$ est fermé et si $G\longrightarrow S$ est
universellement ouvert en les points de $G^\sigma$ (cf. corollaire 1.5
ci-dessous), alors $G^\sigma\longrightarrow S$ est universellement ouvert.

\noindent (iv) Si $G^\circ$ est fermé, il en est de même de $G^\rho$.
Supposons qu'on soit en égale caractéristique, et que, pour tout entier $n>0$
tel que $(n,p)=1$, l'homomorphisme de puissance $n$-ième dans $G$ soit ouvert,
alors $G^\rho$ est ouvert.

Donnons des indications sur la démonstration. Le fait que $G^\circ$ soit
localement constructible, et quasi-compact sur $S$, est contenu dans le lemme
suivant :

\result{LEMME}{1.2} Supposons que $S\neq\varnothing$, alors il existe un ouvert
non vide $U$ dans $S$, et un schéma en groupes de type fini $H$ sur $U$ à
fibres connexes, enfin un homomorphisme de schémas en groupes
$H\longrightarrow G|U$ qui soit une immersion ouverte ayant pour image
$G^\circ|U$.

Pour prouver le lemme, on peut supposer $S$ irréductible, soit $\eta$ son point
générique, et faisons le changement de base
$S'=\operatorname{Spec}(\mathcal O_{S,\eta})\longrightarrow S$ ; on trouve un
schéma en groupes $G'$ sur un anneau artinien local
$\mathcal O_{S,\eta}=A$, dans lequel on a un sous-schéma en groupes ouvert
$G'{}^\circ$ de type fini sur $A$, comme on a dit plus haut. Il provient donc
d'un schéma en groupes de type fini $H$ sur un voisinage ouvert $U$ de $\eta$,
et l'immersion canonique $G'{}^\circ\longrightarrow G'$ provient d'une
immersion

\src{225}
ouverte $H\longrightarrow G|U$, qui sera un homomorphisme de schémas en groupes
pour $U$ assez petit. Comme les fibres de $H$ sont connexes si on prend $U$
assez petit, et toutes de même dimension, égale à celle des fibres de $G$ pour
$U$ assez petit, il s'ensuit que, pour tout $s\in U$, l'image de $H_s$ dans
$G_s$ est exactement $G_s^\circ$ ($U$ assez petit), ce qui prouve 1.2.

La deuxième assertion dans (i) est contenue dans le lemme suivant (qu'on
applique à un voisinage ouvert quasi-compact de $G^\circ$ dans $G$) :

\result{LEMME}{1.3} Soient $X$ un préschéma de type fini et séparé sur $S$
localement noethérien, $g$ une section de $X$ sur $S$, $X^\circ$ la réunion des
composantes connexes des $g(s)$ dans les $X_s$. Soit $s\in S$, tel que
$X_s^\circ$ soit propre sur $k(s)$, alors il existe un voisinage ouvert $U$ de
$s$, tel que $X^\circ|U$ soit propre sur $U$, et a fortiori fermé dans $X|U$.

Par descente fidèlement plate de la base, on se ramène au cas où $S$ est le
spectre d'un anneau local complet, et $s$ son point fermé. Appliquant [2], III,
5.5.1, on voit que $X$ se décompose en somme de deux ouverts disjoints $X'$ et
$X''$, le premier propre sur $S$ et tel que $X'_s=X_s^\circ$. Cela nous ramène
au cas où $X=X'$, i.e. où $X$ est propre sur $S$. Dans ce cas on s'en tire par
une démonstration standard, utilisant le critère valuatif de propreté d'une
partie (oublié dans [2], chap. II).

Prouvons que $G^\tau$ est ouvert, ou ce qui revient au même, compte tenu que la
formation de $G^\tau$ (comme celle de $G^\circ$, $G^\sigma$, $G^\rho$) commute
à l'extension de la base : pour toute section $g$ de $G$ sur $S$,
$g^{-1}(G^\tau)$ est ouvert. Cela signifie deux choses :

\medskip
\noindent a. Soit $y\in S$ tel que $g(y)\in G_y^\tau$, alors pour tout
$y'\in\overline{\{y\}}$ voisin de $y$, on a $g(y')\in G_{y'}^\tau$.

\medskip
\noindent b. Soit $y'\in S$ tel que $g(y')\in G_{y'}^\tau$, alors, pour toute
générisation $y$ de $y'$, on a $g(y)\in G_y^\tau$.

Pour (a), on note qu'il existe un entier $n>0$ tel que
$g^n(y)\in G_y^\circ$ ; comme $G^\circ$ est constructible,
$(g^n)^{-1}(G^\circ)$ est constructible, il s'ensuit que l'on a
$g^n(y')\in G_{y'}^\circ$ pour $y'\in\overline{\{y\}}$ voisin de $y$. Pour
(b), on note que les $g(y')^n=g^n(y')$ restent dans un ouvert quasi-compact de
$G_{y'}$ (car ils sont contenus dans un nombre fini de classes modulo
$G_{y'}^\circ$), donc il existe un ouvert quasi-compact $U$ dans $G$ qui
contient les $g^n(y')$, donc aussi leurs générisations $g^n(y)$, donc les
puissances de $g(y)$ restent dans un ouvert quasi-compact de $G_y$, ce qui
implique facilement que $g(y)\in G_y^\tau$.

\src{226}
Supposons $G^\circ$ fermé, prouvons qu'il en est de même de $G^\tau$. Comme on
sait déjà que $G^\tau$ est ouvert, donc localement constructible, il reste à
prouver qu'il est stable par spécialisation, ce qui provient du fait que c'est
une réunion de fermés, à savoir les images inverses, par les homomorphismes
$\varphi_n$ de puissance $n$-ième, du fermé $G^\circ$.

Le même argument prouvera que $G^\sigma$ et $G^\rho$ sont fermés si $G^\circ$
l'est (sous réserve dans le premier cas qu'on soit en égale caractéristique),
une fois démontré que $G^\sigma$ et $G^\rho$ sont localement constructibles.
Or, pour $x\in G^\tau$, soit $\nu(x)$ le plus petit entier $n>0$ tel que
l'homomorphisme $\varphi_n$ de puissance $n$-ième envoie $x$ dans $G^\circ$.
Donc $G^\sigma$ (resp. $G^\rho$) est formé des $x\in G$\footnote{Le texte
imprimé porte ici « $X$ », qui n'est pas défini dans ce passage ; les quatre
lieux considérés sont des parties de $G$.} tels que $\nu(x)$ soit premier à
$p$ (resp. une puissance de $p$), et notre assertion de constructibilité
résulte alors de la suivante, plus précise :

\result{LEMME}{1.4} La fonction $\nu$ sur $G^\tau$ est localement
constructible.

Cela signifie en effet que, pour tout entier $n>0$, l'ensemble des
$x\in G^\tau$ tels que $\nu(x)=n$ est localement constructible ; or c'est la
différence de $\varphi_n^{-1}(G^\circ)$ et de la réunion des
$\varphi_d^{-1}(G^\circ)$, où $d$ parcourt les diviseurs propres de $n$ ; comme
$G^\circ$ est localement constructible, il en est de même des
$\varphi_d^{-1}(G^\circ)$, donc aussi de la différence précédente.

Supposons $G^\circ$ fermé et $G\longrightarrow S$ universellement ouvert en les
points de $G^\sigma$, prouvons que $G^\sigma\longrightarrow S$ est ouverte,
i.e. transforme un voisinage dans $G^\sigma$ d'un $x\in G^\sigma$ en un
voisinage de $y=f(x)$. Comme $G^\sigma$ est localement constructible, on est
ramené à prouver que, pour toute générisation $y'$ de $y$, il existe une
générisation $x'$ de $x$ dans $G^\sigma$ au-dessus de $y'$. Cela nous ramène
par changement de base au cas où $S$ est le spectre d'un anneau de valuation
discrète et où $y,y'$ sont respectivement le point fermé et le point générique.
Utilisant le fait que $G\longrightarrow S$ est ouvert en $x$ (donc qu'il existe
une générisation $x_1$ de $x$ dans $G$ au-dessus de $y'$), on peut supposer de
plus qu'il existe une section $g$ de $G$ sur $S$, telle que $x=g(y)$ (quitte à
faire encore un changement de base). Si $k(y')$ est de caractéristique 0, il
suffit de prendre n'importe quelle générisation $x'$ de $x$ dans $G$ au-dessus
de $y'$ ; elle est dans $G^\tau$ puisque $G^\tau$ est ouvert, donc dans
$G^\sigma$ puisque $G_{y'}^\sigma=G_{y'}^\tau$. Si la caractéristique de
$k(y')$ est $p>0$, posons
\[
\nu(g(y'))=p^hm\qquad\text{avec }(m,p)=1 ;
\]

\src{227}
soient $a$ et $b$ deux entiers tels que $ap^h+bm=1$, posons
\[
g_1=g^{ap^h},\qquad g_2=g^{bm},\qquad\text{d'où}\qquad g=g_1g_2 ;
\]
donc, par construction, on a $g_1(y')\in G^\sigma$ et
$g_2(y')\in G^\rho$. Comme $G^\circ$ est fermé, il s'ensuit que
$g_1(y)\in G^\sigma$ et $g_2(y)\in G^\rho$, d'où, en vertu de
\[
g(y)=g_1(y)g_2(y)\in G^\sigma,
\]
aussi $g_2(y)\in G^\sigma$, donc
\[
g_2(y)\in G_y^\sigma\cap G_y^\rho=G_y^\circ.
\]
Or, de l'hypothèse et du fait que $S$ est le spectre d'un anneau de valuation
discrète, résulte que $G-(G_y-G_y^\circ)$ est un ouvert de $G$ sur lequel
$G\longrightarrow S$ induit un morphisme ouvert, donc que, par tout point de
$G_y^\circ$, passe une « quasi-section » ; par suite, quitte à faire encore une
extension sur la base $S$, on peut supposer qu'il existe une section $g'_2$ de
$G^\circ$ sur $S$ telle que $g'_2(y)=g_2(y)$. Posons $g'=g_1g'_2$ ; alors, par
construction, $g'(y)=g(y)=x$, et
\[
g'(y')=g_1(y')g'_2(y')\in G^\sigma,
\]
donc $g'(y')=x'$\footnote{Le texte imprimé porte ici « $g(y')=x'$ » ; la
section construite dans la phrase précédente et dont on calcule la fibre
générique est $g'$.} est une générisation de $x$ dans $G^\sigma$ au-dessus de
$y'$, ce qui prouve (iii).

Prouvons enfin la dernière assertion (iv). On est ramené à prouver que, si
$x\in G^\rho$, alors toute générisation $x'$ de $x$ est dans $G^\rho$. On peut
supposer (quitte à prendre les images par $\varphi_{p^h}$ pour $h$
convenable)\footnote{Le texte imprimé porte $\varphi_{n^h}$ ; pour
$x\in G^\rho$, la définition donnée p.~223 fournit une puissance de l'exposant
caractéristique $p$ qui envoie $x$ dans $G^\circ$.} que l'on a même
$x\in G^\circ$. Alors, pour tout entier $n$ premier à la caractéristique, $x$
est dans l'image de $\varphi_n$ (car la puissance $n$-ième dans un groupe de
type fini connexe sur un corps de caractéristique première à $n$ est
surjective). Comme $\varphi_n$ est ouverte, il s'ensuit que $x'$ est également
dans l'image de $\varphi_n$. Plus précisément, soit $U$ un ouvert
quasi-compact de $G^\tau$ contenant $G_y^\circ$, alors on a
$x'\in\varphi_n(U)$ pour tout $n$ premier à $p$. Prenant pour $n$ un multiple
commun des facteurs premiers à $p$ des $\nu(z)$, pour $z\in U$, on trouve que
$x'\in G^\rho$.

L'assertion (1.1), (iii), se complète ainsi :

\result{COROLLAIRE}{1.5} Soit $n>1$ un entier tel que l'homomorphisme de
puissance $n$-ième $\varphi_n:G\longrightarrow G$ soit universellement ouvert
(par exemple étale), posons
\[
G^{(n)}=\varphi_n^{-1}(G^\circ),
\]
et supposons que, dans chaque fibre géométrique $G_{\bar s}$, le sous-ensemble

\src{228}
$\bigcup_{h\geq0}|{}_{n^h}G_{\bar s}|$ soit dense dans
$\bigcup_{h\geq0}G_{\bar s}^{(n^h)}$. Alors $G\longrightarrow S$ est universellement ouvert en les
points des $G^{(n^h)}$. En particulier, si les $G_s^\circ$ ne contiennent pas
de composante additive et si, pour tout entier $n>1$, l'homomorphisme
$\varphi_n$ est universellement ouvert en les points de caractéristique
résiduelle première à $n$\footnote{Le texte imprimé porte « première à $p$ ».
Dans cette proposition quantifiée par $n$, le critère d'étalité et l'argument
par les noyaux ${}_nG$ exigent une caractéristique résiduelle première à $n$.},
alors $G\longrightarrow S$ est universellement ouvert en les points de
$G^\sigma$, donc, par (1.1), (iii), si $G^\circ$ est fermé,
$G^\sigma\longrightarrow S$ est universellement ouvert. Dans ces cas,
$s\longmapsto\dim G_s$ est une fonction localement constante sur $S$.

En effet, de l'hypothèse résulte que le noyau ${}_nG$ de $\varphi_n$ est
universellement ouvert sur $S$, donc $G\longrightarrow S$ est universellement
ouvert en les points de ${}_nG$, donc (remplaçant $n$ par $n^h$) en les points
des ${}_{n^h}G$. L'hypothèse de densité dit exactement que, fibre par fibre, les points d'ordre
une puissance de $n$ sont denses dans la réunion des $G^{(n^h)}$. Pour
l'assertion particulière, l'absence de composantes additives implique cette
densité pour les entiers $n$ inversibles dans le corps résiduel ; d'où résulte facilement que
$G\longrightarrow S$ est universellement ouvert en tous les points de cette
réunion, en particulier le long de la section unité, d'où facilement le fait
que $s\longmapsto\dim G_s$ est une fonction localement constante.

\result{REMARQUE}{1.6} Rappelons qu'un morphisme $X\longrightarrow S$ est dit
universellement ouvert s'il transforme tout ouvert en un ouvert, et garde cette
propriété louable après tout changement de base $S'\longrightarrow S$. Cela
signifie aussi (lorsque $S$ est localement noethérien et $X\longrightarrow S$
localement de type fini) que toute composante irréductible de $X$ domine $S$,
et que cette propriété est conservée après tout changement de base
$S'\longrightarrow S$. Dans ces deux assertions, il suffit d'ailleurs
(moyennant les hypothèses de finitude ci-dessus) de tester avec des changements
de base $S'\longrightarrow S$, où $S'$ est le spectre d'un anneau de valuation
discrète (complet à corps résiduel algébriquement clos si on y tient~...). La
définition s'étend de façon évidente au cas d'une partie $Z$ de $X$ (telle la
partie $G^\sigma$ de $G$). C'est un phénomène parfaitement normal, même si on
part d'un morphisme projectif simple $X\longrightarrow S$ à fibres
géométriques connexes (par exemple le carré fibré de la famille modulaire de
courbes elliptiques sur $S=\operatorname{Spec}\mathbf C[j]$), que
$\operatorname{Pic}_{X/S}$ ne soit pas universellement ouvert sur $S$, i.e.
qu'il puisse y avoir des composantes irréductibles de
$\operatorname{Pic}_{X/S}$ se trouvant tout entières au-dessus d'un seul point
de $S$ : cela est lié au fait que le rang du groupe de Néron--Severi des fibres
de $X/S$ peut faire des sauts vers le haut (phénomènes de « multiplication
complexe »). Par contre, (1.5) nous assure que, dans les bons cas,
$\operatorname{Pic}_{X/S}^\sigma$ (et le plus souvent même, semble-t-il,
$\operatorname{Pic}_{X/S}^\tau$) est universellement ouvert sur $S$.

\src{229}
Enfin, voici un cas utile où, exceptionnellement, $G^\circ$ consente à être
ouvert :

\result{COROLLAIRE}{1.7} Supposons que $G\longrightarrow S$ soit
universellement ouvert en les points de $G^\circ$ (cf. (1.5)), et que les
fibres $G_s$ soient séparables, donc simples sur $k(s)$ (cette dernière
condition étant automatiquement remplie en caractéristique résiduelle nulle,
en vertu d'un résultat de Cartier). Alors $G^\circ$ est ouvert dans $G$. Si de
plus $S$ est réduit, $G^\circ$ est simple (en particulier plat) sur $S$.

La première assertion peut d'ailleurs se préciser en notant que si, pour un
$s\in S$ donné, $G\longrightarrow S$ est universellement ouvert en les points
de $G_s^\circ$ et si $G_s^\circ$ est séparable sur $k(s)$, alors $G^\circ$ est
un voisinage de $G_s^\circ$ (et d'ailleurs, l'hypothèse faite en $s$ restera
vérifiée en les points voisins). Cet énoncé, d'ailleurs indépendant de toute
structure de groupes, se trouvera démontré dans [2], IV, paragraphe 7. La
dernière assertion dans (1.7), également indépendante de toute structure de
groupe et de toute question de composante connexe, est un cas particulier d'un
critère de platitude donné dans [2], IV, paragraphe 5, qui implique, plus
généralement :

\result{COROLLAIRE}{1.8} Soit $U$ une partie ouverte d'une fibre $G_s$, telle
que $G\longrightarrow S$ soit universellement ouvert en les points de $U$ (cf.
(1.5)). Si $G_s$ est séparable sur $k(s)$ et $S$ réduit en $s$, alors $G$ est
plat sur $S$ en les points de $U$ (donc en l'occurrence, simple sur $S$ en les
points de $U$).

\result{REMARQUES}{1.9} J'ignore si, lorsque $G$ est séparé sur $S$, $G^\circ$
ou $G^\tau$ est toujours fermé dans $G$ ; cela semble peu probable. On trouve
en tous cas des contre-exemples évidents si on laisse tomber l'hypothèse de
séparation, par exemple en prenant la droite affine dont on dédouble une
infinité dénombrable de fois l'origine, obtenant ainsi un schéma en groupes $G$
sur la droite affine dont toutes les fibres sont réduites au groupe unité, sauf
une dont la fibre est $\mathbf Z$. Ce schéma en groupes est un sous-schéma en
groupes ouvert, savoir l'adhérence de la section unité, dans le préschéma de
Picard du $S$-schéma $X$ correspondant à une famille de coniques dégénérant en
deux droites concourantes.

D'ailleurs, même en partant d'un schéma en groupes fini et plat $G$ sur le
spectre $S$ d'un anneau de valuation discrète $V$, et par suite si $G^\circ$
est réduit à la section unité, donc fermé, diverses conclusions deviennent
fausses si on abandonne certaines hypothèses. Supposons que $V$ soit d'égale
caractéristique $p>0$, et soit $G$ le noyau de l'homomorphisme
$\mathbf G_a\longrightarrow\mathbf G_a$ défini par l'homomorphisme de
foncteurs
\[
f\longmapsto f^p-tf,
\]
où $t$ est une uniformisante de $V$.

\src{230}
(Se rappeler que, par définition, le « groupe additif » $\mathbf G_a$ sur $S$
représente le foncteur $S'\longmapsto\Gamma(S',\mathcal O_{S'})$.) Alors, en
tant que $S$-schéma,
\[
G=\operatorname{Spec}V[x]/(x^p-tx)
\]
est réunion de $p$ « droites » concourantes, figurant un groupe
$\mathbf Z/p\mathbf Z$ qui dégénère en un groupe infinitésimal de type additif.
Dans cet exemple, $G^\circ=G^\sigma$ (et c'est l'une des $p$ droites), donc
n'est pas ouvert dans $G$. Dans le cas d'inégale caractéristique, de
caractéristique résiduelle $p>0$, partons du schéma en groupes
$H=\boldsymbol\mu_p$, noyau de la puissance $p$-ième dans $\mathbf G_m$, donc
\[
H=\operatorname{Spec}V[x]/(x^p-1) ;
\]
c'est encore la réunion de $p$ droites concourantes, figurant un groupe
$\mathbf Z/p\mathbf Z$ (en caractéristique 0) qui dégénère en un groupe
infinitésimal de type multiplicatif en caractéristique $p$. Soit $H'$ le
« schéma en groupes constant » défini par le groupe fini ordinaire
$\mathbf Z/p\mathbf Z$, donc la somme disjointe de $\mathbf Z/p\mathbf Z$
copies de $S$, ou encore $\operatorname{Spec}V^{\mathbf Z/p\mathbf Z}$. Alors
$G=H\times_S H'$ décrit un groupe $(\mathbf Z/p\mathbf Z)^2$ en
caractéristique 0, dégénérant en un groupe infinitésimal fois
$\mathbf Z/p\mathbf Z$ en caractéristique $p$. Ici $G^\rho$ est la réunion de
la section unité et de la fibre spéciale, donc $G^\rho$ n'est pas ouvert,
alors que $G^\sigma$ est la réunion de la section unité et de la fibre
générale, donc n'est pas fermé, contrairement à ce qui est affirmé dans les
cas d'égale caractéristique dans (1.1), (iii) et (iv). Bien entendu, ce sont là
des phénomènes liés à la caractéristique $p>0$. Les résultats qui précèdent
donnent en effet :

\result{COROLLAIRE}{1.10} Supposons les caractéristiques résiduelles de $S$
toutes nulles, donc $G^\tau=G^\sigma$, et $G^\circ=G^\rho$. Alors
$G^\tau=G^\sigma$ est ouvert, et même ouvert et fermé si $G$ est séparé sur
$S$ et si les $G_s^\circ$ sont propres ; sous cette même hypothèse,
$G^\circ=G^\rho$ est propre sur $S$, donc fermé dans $G$. Supposons enfin que,
pour tout entier $n>1$, l'homomorphisme de puissance $n$-ième dans $G$ soit
universellement ouvert, alors $G^\circ$ est ouvert ; et si de plus les
$G_s^\circ$ n'ont pas de composante additive (par exemple les $G_s^\circ$
propres, comme ci-dessus), alors $G^\tau\longrightarrow S$ est universellement
ouvert, et même simple si $S$ est réduit.

Signalons enfin le résultat facile suivant :

\result{PROPOSITION}{1.11} Il existe une partie ouverte $U$ de $S$, telle que
l'ensemble $G'$ des points de $G$ en lesquels $G$ est simple (resp. plat)
au-dessus de $S$ soit l'ensemble sous-jacent à un sous-schéma en groupes ouvert
induit de $G|U$. De plus, toute section de $G$ sur $U$ est une section de $G'$
sur $U$.

\src{231}
\result{COROLLAIRE}{1.12} Si $G$ est simple (resp. plat) sur $S$ en les points
de la section unité, elle l'est en les points de toute section de $G$ sur $S$,
et en tous les points de $G^\circ$. Si de plus, pour tout entier $n>0$,
l'homomorphisme de puissance $n$-ième
$\varphi_n:G\longrightarrow G$ est étale en les points de caractéristique
première à $n$, alors $G$ est simple (resp. plat) sur $S$ en tous les points de
$G^\sigma$.

\fsection{2}{Application aux propriétés locales des schémas de Picard}

\result{THÉORÈME}{2.1}

\noindent(i) Soit $f:X\longrightarrow S$ un morphisme propre et simple, et
supposons que $\operatorname{Pic}_{X/S}$ existe (par exemple $f$ un morphisme
projectif). Alors $\operatorname{Pic}_{X/S}$ est séparé sur $S$, et, pour toute
partie fermée $Z$ de $\operatorname{Pic}_{X/S}$ qui est de type fini sur $S$,
$Z$ est propre sur $S$.

\noindent(ii) Soit $X$ un préschéma sur un corps $k$, propre et géométriquement
normal. Alors $\operatorname{Pic}_{X/k}^\circ$\footnote{Le texte imprimé porte
$\operatorname{Pic}_{X/S}^\circ$, alors que $S$ n'intervient pas dans cet
énoncé sur le corps $k$.} est propre sur $k$.

\begin{proof}
(i) Les critères valuatifs ([2], II, paragraphe 7) nous ramènent à prouver ceci :
si $S$ est le spectre d'un anneau de valuation discrète complet, $U$ l'ouvert
réduit au point générique de $S$, alors toute section rationnelle de
$\operatorname{Pic}_{X/S}$ sur $S$, i.e. toute section au-dessus de $U$, se
prolonge de façon unique en une section au-dessus de $S$. Compte tenu de la
définition de $\operatorname{Pic}_{X/S}$, cela équivaut à l'énoncé suivant :
pour tout Module inversible $\mathcal L$ sur $V=f^{-1}(U)$, il existe un Module
inversible sur $X$ qui prolonge $\mathcal L$, unique à un isomorphisme près. Or
ceci résulte facilement de la description des Modules inversibles sur $V$,
resp. $X$, en termes de classes de diviseurs « de Cartier », compte tenu que
les anneaux locaux de $X$ sont réguliers ($X$ étant simple sur $S$ régulier),
donc factoriels en vertu d'Auslander, ce qui implique que tout diviseur sur
$X$\footnote{Le texte imprimé porte ici « sur $S$ » ; la factorialité invoquée
est celle des anneaux locaux de $X$.} est un diviseur de Cartier. En effet,
tout diviseur sur $V$ peut se prolonger en un diviseur sur $X$ en prenant son
« adhérence ».
\end{proof}

\result{REMARQUE}{2.2} La démonstration reste valable en supposant seulement
que $f$ est plat et ses fibres $X_s$ sont localement des intersections
complètes et simples sur $k(s)$ en codimension $\leq2$, compte tenu du résultat
suivant démontré dans [5] : un anneau local noethérien intersection complète
qui est régulier en codimension $\leq3$ est factoriel (« conjecture de Samuel »).
On notera que le

\src{232}
résultat devient faux si on remplace « codimension $\leq2$ » par « codimension
$\leq1$ », i.e. par l'hypothèse « normal », comme on peut se convaincre sur
l'exemple d'une famille de quadriques non singulières dégénérant en un cône
quadratique.

(ii) Utilisant le lemme de Chow, on est ramené au cas où $X$ est projectif,
donc plongé dans $\mathbf P_k^n$ ; on peut supposer $X$ connexe. Si
$\dim X=1$, alors $X$ est simple sur $k$, et on applique (i). Si $\dim X\geq2$,
on peut utiliser une variante des « critères d'équivalence » connus, qui
implique qu'il existe un nombre fini de courbes $Y_i$ simples sur $X$ (obtenues
comme intersections de $X$ avec des sous-espaces linéaires convenables de
$\mathbf P_k^n$), telles que
\[
\operatorname{Pic}_{X/k}^\tau\longrightarrow
\prod_i\operatorname{Pic}_{Y_i/k}^\tau
\]
soit un monomorphisme, et induit a fortiori un monomorphisme pour les
composantes connexes. Comme celle du second membre est propre sur $k$ d'après
ce qui précède, et qu'il s'agit d'un homomorphisme de schémas en groupes, qui
sera nécessairement une immersion fermée, il s'ensuit que
$\operatorname{Pic}_{X/k}^\circ$ est également propre sur $k$. On peut éviter
le recours aux délicats critères d'équivalence en utilisant la structure des
groupes algébriques commutatifs sur un corps algébriquement clos (due à
Chevalley--Borel) ; on est ramené à prouver que tout morphisme de la droite
affine privée de l'origine\footnote{Le texte imprimé porte « de la droite
affine ». L'erratum lié, p.~303, prescrit « droite affine privée de
l'origine ».} dans $\operatorname{Pic}_{X/k}^\tau$ est constant, ce qui
équivaut à dire que tout Module inversible sur
$X[t,t^{-1}]$\footnote{Le texte imprimé porte $X[t]$ ; l'erratum lié, p.~303,
prescrit $X[t,t^{-1}]$.} provient d'un Module inversible sur $X$, résultat qui
est bien connu et fort élémentaire et n'utilise pas même le fait que $X$ soit
propre sur $k$ (l'hypothèse $X$ normal permettant de se réduire aussitôt au cas
$X$ régulier).

\result{COROLLAIRE}{2.3} Soit $f:X\longrightarrow S$ un morphisme propre et
normal (i.e. plat et à fibres géométriques normales), supposons que
$\operatorname{Pic}_{X/S}$ existe, alors $\operatorname{Pic}_{X/S}^\circ$ est
propre sur $S$, donc fermé dans $\operatorname{Pic}_{X/S}$ ; de plus
$\operatorname{Pic}_{X/S}^\tau$ et $\operatorname{Pic}_{X/S}^\rho$ sont
également fermés, et aussi $\operatorname{Pic}_{X/S}^\sigma$ en « égale
caractéristique ».

On applique (1.1) et (2.1), (ii).

\result{COROLLAIRE}{2.4} Soit $f:X\longrightarrow S$ un morphisme propre et
simple, tel que $\operatorname{Pic}_{X/S}$ existe et soit somme de schémas
$P^{(i)}$ de type fini sur $S$ (cf. [3], V, 4.1). Alors chaque $P^{(i)}$ est
propre sur $S$.

Résulte de (2.1), (i). Comme on l'a signalé dans (2.2), le résultat peut se
généraliser en faisant des hypothèses moins restrictives sur les fibres de $f$,
mais devient faux si on se borne à supposer $f$ normal. Dans ce cas, j'ignore
d'ailleurs si $\operatorname{Pic}_{X/S}^\tau$ est néanmoins propre sur $S$,
même en admettant qu'il soit de type fini sur $S$.

\src{233}
\result{THÉORÈME}{2.5} Soit $f:X\longrightarrow S$ un morphisme propre et
plat, tel que $\operatorname{Pic}_{X/S}$ existe, et pour tout entier $n$, soit
$\varphi_n$ l'homomorphisme puissance $n$-ième dans ce préschéma en groupes.
Alors $\varphi_n$ est étale en tous les points
$x\in\operatorname{Pic}_{X/S}$\footnote{Le texte imprimé porte « $x\in X$ » ;
$\varphi_n$ est l'endomorphisme de puissance du préschéma de Picard.} à
caractéristique résiduelle première à $n$.

Par la caractérisation infinitésimale des morphismes étales, cette assertion
équivaut à la suivante :

\result{LEMME}{2.6} Supposons que $S$ soit le spectre d'un anneau local
artinien $A$ dont l'idéal maximal $\mathfrak m$ soit de puissance
$(\nu+1)$-ième nulle, soit
\[
A_{\nu-1}=A/\mathfrak m^\nu,
\qquad X_{\nu-1}=X\otimes_A A_{\nu-1},
\]
$\mathcal L$ un Module inversible sur $X$, enfin $\mathcal L'_{\nu-1}$ un
Module inversible sur $X_{\nu-1}$ dont la puissance tensorielle $n$-ième soit
isomorphe à
\[
\mathcal L_{\nu-1}=\mathcal L\otimes_A A_{\nu-1}.
\]
Alors il existe un Module inversible $\mathcal L'$ sur $X$ dont la puissance
tensorielle $n$-ième est isomorphe à $\mathcal L$ (si $n$ est premier à la
caractéristique résiduelle de $k=k(A)$).

Prouvons le lemme. On pose
\[
V=\mathfrak m^\nu=\mathfrak m^\nu/\mathfrak m^{\nu+1} ;
\]
c'est un espace vectoriel sur $k=k(A)$. On commence à prolonger
$\mathcal L'_{\nu-1}$ en un Module inversible quelconque $\mathcal L'$ sur $X$.
L'obstruction à ce faire se trouve dans
\[
H^2(X_0,\mathcal O_{X_0})\otimes_k V,
\]
mais, en vertu de l'hypothèse faite sur $\mathcal L'_{\nu-1}$ et du fait que
$\mathcal L_{\nu-1}$ peut se prolonger, on voit que le produit de cette
obstruction par $n$ est nul, donc elle-même est nulle puisque $n$ est premier
à la caractéristique. D'ailleurs, l'arbitraire dans le prolongement effectué
se trouve dans
\[
H^1(X_0,\mathcal O_{X_0})\otimes_k V,
\]
et la déviation $\xi$ de $\mathcal L'^{\otimes n}$ avec $\mathcal L$ se trouve
dans le même Module ; si on essaye de corriger $\mathcal L'$ de façon à rendre
cette déviation nulle, on est ramené à trouver un $\eta$ dans ledit Module,
tel que $n\eta=\xi$. Or c'est possible encore grâce au fait que $n$ est
premier à la caractéristique.

\result{COROLLAIRE}{} Sous les conditions de (2.5), supposons de plus que les
schémas de Picard des fibres $X_s$ ne contiennent pas de composante additive
(par exemple les $X_s$ géométriquement normaux, cf. (2.1), (ii)). Alors
$\operatorname{Pic}_{X/S}\longrightarrow S$ est universellement ouvert en les
points de $\operatorname{Pic}_{X/S}^\circ$. Si
$\operatorname{Pic}_{X/S}^\circ$ est fermé (par exemple les $X_s$
géométriquement normaux, cf. (2.3)), alors
$\operatorname{Pic}_{X/S}^\sigma$ est lui-même universellement ouvert sur $S$.
Enfin, dans le cas d'égale caractéristique,
$\operatorname{Pic}_{X/S}^\rho\longrightarrow S$ est universellement ouvert.

On applique (1.5) et (1.1).

\src{234}
\result{COROLLAIRE}{2.7} Soit $f:X\longrightarrow S$\footnote{Le texte
imprimé porte « $f:X\longrightarrow Y$ », alors que le préschéma de Picard et
la fonction qui suivent ont tous deux pour base $S$.} un morphisme propre et
plat, tel que $\operatorname{Pic}_{X/S}$ existe. Alors la fonction
\[
s\longmapsto\dim\operatorname{Pic}_{X_s/k(s)}
\]
sur $S$ est semi-continue supérieurement (i.e. elle peut faire des sauts vers
le haut, mais pas vers le bas), et elle est même continue (i.e. localement
constante) si les $\operatorname{Pic}_{X_s/k(s)}$ ne contiennent pas de
composante additive.

La première assertion est vraie trivialement ou presque pour tout préschéma en
groupes localement de type fini sur une base localement noethérienne, puisqu'il
suffit de regarder le long de la section unité. La deuxième assertion résulte
de (2.5).

\result{REMARQUE}{2.8} Soient $s,s'\in S$ tels que $s$ soit une spécialisation
de $s'$, alors (2.7) équivaut à une inégalité (resp. égalité) entre les
dimensions des schémas de Picard de $X_{s'}$ et de sa « spécialisée » $X_s$.
D'ailleurs, Serre avait fait observer, dès avant la construction des schémas de
Picard, que l'invariance de la dimension des Picards des $X_s$ dans le cas d'un
morphisme simple $f:X\longrightarrow S$ était une conséquence formelle de la
théorie de la spécialisation du groupe fondamental ([4], X), et des relations
classiques à la Kummer entre les points d'ordre fini sur la variété de Picard,
et le groupe fondamental rendu abélien ([4], XI). Si on désigne de façon
générale par $\alpha,\mu,\lambda$ les dimensions de la partie abélienne,
multiplicative, additive de $\operatorname{Pic}_{X_s/k(s)}^\circ$, et qu'on
définit de même $\alpha',\mu',\lambda'$, les relations connues s'expriment par
les inégalités suivantes :
\[
\tag{*}\alpha+\mu+\lambda\geq\alpha'+\mu'+\lambda'
\]
(valable pourvu que $\operatorname{Pic}_{X/S}$ existe, donc probablement en
tous cas), inégalité qui, pour $\lambda=\lambda'=0$, se réduit à une égalité,
valable sous les mêmes hypothèses d'existence :
\[
\alpha+\mu=\alpha'+\mu' ;
\]
d'autre part on a
\[
\tag{**}2\alpha+\mu\leq2\alpha'+\mu'
\]

\src{235}
par l'argument de Serre, si les $X_s$ sont séparables (sans même supposer
l'existence de $\operatorname{Pic}_{X/S}^\tau$), ou si les
${}_n\operatorname{Pic}_{X/S}$ (noyaux de $\varphi_n$ dans le foncteur de
Picard) sont séparés sur l'ouvert de $S$ où $n$ est inversible (compte tenu
qu'ils sont étales sur cet ouvert grâce à (2.5)). On aurait tendance à conjecturer que (*) est une égalité en tous cas,
ou du moins si les $X_s$ sont séparables, et qu'on a des inégalités
\[
\tag{***}\alpha\leq\alpha',\qquad\lambda\geq\lambda',
\]
qui devraient être valables chaque fois qu'on a un préschéma en groupes,
localement de type fini sur $S$ localement noethérien, dans lequel la dimension
des fibres est constante (voir (1.3) pour un résultat positif dans cette
direction).

\result{REMARQUE}{2.9} Dans tous les cas connus,
$\operatorname{Pic}_{X/S}^\tau$ est universellement ouvert sur $S$, mais il ne
faut sans doute pas s'en autoriser pour des illusions excessives, même si
$f:X\longrightarrow S$ est simple ; en tous cas, Mumford a construit un
exemple (il est vrai avec $S$ non réduit, en fait $S$ spectre d'un anneau
artinien) où $\operatorname{Pic}_{X/S}^\tau$ n'est pas plat sur $S$, en faisant
varier infinitésimalement la surface de Igusa. Le point envisagé par Mumford se
trouve d'ailleurs dans $\operatorname{Pic}_{X/S}^\rho$, et il reste possible
(pour $f:X\longrightarrow S$ simple) que $\operatorname{Pic}_{X/S}^\tau$ soit
plat sur $S$ aux points de $\operatorname{Pic}_{X/S}^\sigma$ ; le conférencier
doute cependant qu'il en soit toujours ainsi, même en se restreignant aux
points de $\operatorname{Pic}_{X/S}^\circ$ et en supposant que $S$ est le
spectre d'un anneau de valuation discrète. La question est d'ailleurs liée à
l'étude des points fixes d'un schéma abélien sous un groupe fini
d'automorphismes, situation pour laquelle on semble manquer d'exemples.

Il semblerait que, même en se bornant à $f:X\longrightarrow S$ simple et
projectif, les résultats de régularité locale sur $\operatorname{Pic}_{X/S}$
énoncés dans le présent numéro, et les conjectures signalées dans (2.8),
épuisent à peu près ce qu'on peut dire à ce sujet sans hypothèses plus
particulières sur la nature des fibres de $f$. Rappelons cependant que, si les
fibres géométriques de $\operatorname{Pic}_{X/S}$ sont réduites et sans
composante additive, alors il résulte de (1.8) et de (2.5) que
$\operatorname{Pic}_{X/S}$ est simple sur $S$ en les points de
$\operatorname{Pic}_{X/S}^\sigma$ lorsque $S$ est réduit ; ce résultat vaut
même, si $f:X\longrightarrow S$ est normal, sans hypothèse sur $S$, comme nous
verrons dans (3.5). Signalons à ce propos :

\result{PROPOSITION}{2.10}

\noindent(i) Si $\operatorname{Pic}_{X/S}$ est simple (resp. plat) sur $S$ en
les points de la section unité, il l'est en tous les points de
$\operatorname{Pic}_{X/S}^\sigma$, et en les points de toute section de
$\operatorname{Pic}_{X/S}$ sur $S$.

\src{236}
\noindent(ii) Soit $s\in S$ tel que
\[
H^2(X_s,\mathcal O_{X_s})=0,
\]
alors il existe un voisinage ouvert $U$ de $s$ tel que
$\operatorname{Pic}_{X/S}|U$ soit simple sur $U$.

\noindent(iii) Soit $X$ un schéma propre sur un corps, alors on a
\[
\dim\operatorname{Pic}_{X/k}\leq\dim H^1(X,\mathcal O_X),
\]
l'égalité ayant lieu si et seulement si $\operatorname{Pic}_{X/k}$ est simple
sur $k$ ; c'est toujours le cas si $k$ est de caractéristique nulle.

(i) résulte de (2.5) et (1.12), (ii) du critère infinitésimal pour les
morphismes simples et d'un calcul d'obstructions bien connu, compte tenu que
(par le « théorème de semi-continuité ») l'hypothèse faite en $s$ sera encore
vérifiée aux points voisins. Enfin (iii) résulte du fait que
$H^1(X,\mathcal O_X)$ est isomorphe à l'espace tangent de Zariski en l'élément
neutre de $\operatorname{Pic}_{X/k}$ ; la dernière assertion est un cas
particulier d'un théorème de Cartier, disant qu'un « groupe formel » en
caractéristique 0 est formellement simple sur $k$.

\fsection{3}{Le sous-schéma abélien canonique de
$\operatorname{Pic}_{X/S}$, et schéma d'Albanese}

\result{PROPOSITION}{3.1} Soient $k$ un corps, $G$ un schéma en groupes de
type fini sur $k$, commutatif et « sans composante additive ». Alors
$G_{\mathrm{réd}}^\circ$ est séparable sur $k$, et par suite c'est un
sous-schéma en groupes simple sur $k$.

La chose étant triviale si $k$ est parfait, en particulier pour $G_{\bar k}$,
où $\bar k$ est la clôture algébrique de $k$, il suffit de voir que
$(G_{\bar k}^\circ)_{\mathrm{réd}}$ provient d'un sous-schéma de $G$. Or, de
l'hypothèse sur $G_{\bar k}$ (pas de composante additive)\footnote{Le texte
imprimé porte « Or l'hypothèse ... » ; l'erratum lié, p.~303, prescrit « Or,
de l'hypothèse ... ».} résulte facilement
qu'il existe un entier $m$ tel que $(G_{\bar k}^\circ)_{\mathrm{réd}}$ soit
l'image « schématique » de l'homomorphisme de puissance $m$-ième dans
$G_{\bar k}$. Comme ce dernier homomorphisme provient de l'homomorphisme
analogue pour $G$, l'image schématique de ce dernier répond à la question.

\result{COROLLAIRE}{3.2} Soit $X$ un schéma géométriquement normal et propre sur $k$ tel que
$\operatorname{Pic}_{X/k}$ existe, alors il existe un sous-schéma abélien $A$
de $\operatorname{Pic}_{X/k}$ dont l'ensemble sous-jacent est
$\operatorname{Pic}_{X/k}^\circ$.

En effet, en vertu de (2.1), (ii), $\operatorname{Pic}_{X/k}^\circ$ étant
propre sur $k$ satisfait à la condition de (3.1).

\src{237}
Le résultat précédent montre donc que, dans certains cas, la classique
« variété de Picard » (qui serait notée
$(\operatorname{Pic}_{X/k})_{\mathrm{réd}}^\circ$ dans la théorie actuelle)
« est définie sur $k$ », sans supposer le corps $k$ parfait.

Soit maintenant $f:X\longrightarrow S$ un schéma relatif propre et plat, avec
\[
\mathcal O_S\xrightarrow{\sim}f_*(\mathcal O_X)
\]
pour simplifier, tel que $\operatorname{Pic}_{X/S}$ existe et que
$\operatorname{Pic}_{X/S}^\circ$ soit propre sur $S$. Nous supposerons de plus
pour (3.3), (ii), qu'il existe un ouvert de $\operatorname{Pic}_{X/S}$
contenant $\operatorname{Pic}_{X/S}^\circ$ qui soit quasi-projectif sur $S$ ;
cette condition est vérifiée, on l'a vu, si $f$ est projectif et à fibres
géométriques séparables et irréductibles. Rappelons qu'on appelle schéma
abélien sur $S$ un schéma en groupes sur $S$, propre et simple sur $S$, à
fibres géométriques connexes. Nous nous proposons d'examiner s'il existe un
sous-schéma en groupes $A$ de $\operatorname{Pic}_{X/S}$ qui soit un schéma
abélien et dont l'ensemble sous-jacent soit $\operatorname{Pic}_{X/S}^\circ$.
On vient de voir qu'il en existe toujours si $S$ est le spectre d'un corps.
Voici ce qu'on sait dire dans le cas général envisagé ici :

\result{THÉORÈME}{3.3} Sous les conditions précédentes :

\noindent(i) S'il existe un sous-schéma abélien de
$\operatorname{Pic}_{X/S}$ dont l'ensemble sous-jacent soit
$\operatorname{Pic}_{X/S}^\circ$, il est unique. Sa formation est par suite
compatible avec les changements de base.

\noindent(ii) Pour qu'il existe un tel sous-schéma abélien, il suffit qu'il en
existe après tout changement de base $S'\longrightarrow S$, où $S'$ est
artinien local ; si $S$ est le spectre d'un anneau local, il suffit même de
tester avec les $S'=\operatorname{Spec}(A_n)$, où
$A_n=A/\mathfrak m^{n+1}$. Si $S$ est réduit, il suffit également de tester
avec les changements de base $S'\longrightarrow S$, où $S'$ est le spectre
d'un anneau de valuation discrète (complet à corps résiduel algébriquement clos
si on y tient).

\noindent(iii) Supposons que $A$ existe, et soit
$B=\operatorname{Alb}^0(X/S)$ le schéma abélien dual (i.e.
$B=\operatorname{Pic}_{A/S}^\circ$ [7]). Alors on peut construire
canoniquement un espace principal homogène $P=\operatorname{Alb}^1(X/S)$ sous
$B$, et un $S$-morphisme $X\longrightarrow P$ qui soit universel pour les
$S$-morphismes de $X$ dans des schémas para-abéliens (i.e. des espaces
principaux homogènes sous des schémas abéliens). La formation de
$\operatorname{Alb}^0(X/S)$, $\operatorname{Alb}^1(X/S)$ et
$X\longrightarrow\operatorname{Alb}^1(X/S)$ commute au changement de base.

Esquissons la démonstration :

(i) est une propriété de rigidité générale pour les sous-schémas abéliens des
schémas en groupes commutatifs : si deux tels sous-schémas coïncident
ensemblistement en un point $s\in S$, ils coïncident au-dessus de toute la
composante

\src{238}
connexe de $s$ [7]. (Ce résultat généralise un classique théorème de Chow.)

(ii) Utilisant les schémas de Hilbert, on voit que le foncteur qui, à tout
$S'$ sur $S$, associe l'ensemble (réduit à 1 ou 0 éléments) des sous-schémas
abéliens canoniques de $(\operatorname{Pic}_{X/S})\times_S S'$ est
représentable par un schéma de type fini $T$ sur $S$. En vertu de (i),
$T\longrightarrow S$ est un monomorphisme, et en vertu de (3.2) il est
surjectif. Dire qu'il existe un sous-schéma abélien canonique de
$\operatorname{Pic}_{X/S}$ signifie que $T$ a une section sur $S$, ou encore
que $T\longrightarrow S$ est un isomorphisme. Ceci équivaut aussi à dire que
$T\longrightarrow S$ est étale, ou, dans le cas où $S$ est réduit, que
$T\longrightarrow S$ est propre. D'où aussitôt (ii).

(iii) Utilisant simplement la définition de $\operatorname{Pic}_{X/S}$, on
constate que, pour tout schéma abélien $C$ sur $S$, la donnée d'un
$S$-morphisme de $X$ dans un espace principal homogène sous $C$ équivaut à la
donnée d'un homomorphisme de groupes
$C'\longrightarrow\operatorname{Pic}_{X/S}$, où $C'$ est le schéma abélien
dual de $C$. Or, si le sous-schéma abélien canonique $A$ de
$\operatorname{Pic}_{X/S}$ existe, ces homomorphismes se factorisent
nécessairement par $A$ (comme on peut voir encore en utilisant les points
d'ordre fini). D'où aussitôt (iii).

\result{REMARQUES}{3.4} Nous noterons
$\operatorname{Pic}_{X/S}^{\circ\circ}$ le sous-schéma abélien canonique de
$\operatorname{Pic}_{X/S}$, s'il existe. Ce n'est malheureusement pas toujours
le cas, comme on peut voir en faisant varier infinitésimalement la surface de
Igusa (par la déformation modulaire du premier ordre). Il est cependant
possible que $\operatorname{Pic}_{X/S}^{\circ\circ}$ existe du moins si $S$
est réduit, ou, ce qui revient au même d'après (ii), si $S$ est le spectre d'un
anneau de valuation discrète. Soient alors $X_0,X_1$ les fibres spéciale et
générique de $X$, et soit
\[
A_1=\operatorname{Pic}_{X_1/K}^\circ,
\]
où $K$ est le corps des fractions de l'anneau de valuation $V$. En vertu de
Koizumi, il existe un schéma abélien $A$ sur $S$, essentiellement unique, dont
la fibre générale est $A_1$, et on voit aisément comme dans (2.1), (i)
(supposant désormais $X$ simple sur $S$), que le morphisme identique de $A_1$
se prolonge en un morphisme
\[
A\longrightarrow\operatorname{Pic}_{X/S}.
\]
D'où un homomorphisme
\[
\tag{*}A_0\longrightarrow\operatorname{Pic}_{X_0/k}^{\circ\circ},
\]
dont on montre facilement que c'est un homomorphisme surjectif à noyau un
$p$-groupe fini, où $p$ est l'exposant caractéristique du corps résiduel $k$

\src{239}
(toujours par utilisation des points d'ordre fini). Ceci posé, les conditions
suivantes sur $X/S$ sont équivalentes :

\medskip
\noindent a. L'homomorphisme précédent (*) est un isomorphisme (ce que Shimura
exprimerait en disant que la formation de la « variété de Picard » est
« compatible avec spécialisations »).

\noindent b. $\operatorname{Pic}_{X/S}^{\circ\circ}$ existe (alors il ne sera
autre que $A$).

\noindent c. (Pour mémoire) Les
$\operatorname{Pic}_{X_n/S}^{\circ\circ}$ existent.

D'après la remarque qu'on a faite sur le noyau de (*), la condition (a) est
satisfaite si la caractéristique résiduelle est nulle, mais ce résultat sera
généralisé notablement (3.5).

Bien entendu, si $\operatorname{Pic}_{X/S}$ est simple sur $S$ aux points de
$\operatorname{Pic}_{X/S}^\circ$, alors ce dernier est ouvert dans
$\operatorname{Pic}_{X/S}$ (cf. (1.7)) et, par suite, muni de la structure
induite, c'est un sous-schéma abélien de $\operatorname{Pic}_{X/S}$, donc égal
à $\operatorname{Pic}_{X/S}^{\circ\circ}$ qui existe dans ce cas. Mais on a
bien mieux :

\result{THÉORÈME}{3.5} Sous les conditions de (3.3), soit $s\in S$ tel que
$\operatorname{Pic}_{X_s/k(s)}$ soit simple sur $k(s)$ (ou, ce qui revient au
même,
\[
\dim\operatorname{Pic}_{X_s/k(s)}=\dim H^1(X_s,\mathcal O_{X_s})).
\]
Alors il existe un voisinage ouvert $U$ de $s$ tel que
$\operatorname{Pic}_{X/S}$ soit simple sur $S$ aux points de
$\operatorname{Pic}_{X/S}^\circ|U$, qui est donc un sous-schéma abélien ouvert
dans $\operatorname{Pic}_{X/S}|U$. A fortiori,
$\operatorname{Pic}_{X|U/U}^{\circ\circ}$ existe.

Indiquons le principe de la démonstration. Ce qui précède nous ramène au cas où
$S$ est le spectre d'un anneau artinien local $A$, et on raisonne par récurrence
sur l'ordre infinitésimal de $A$. On peut donc supposer
$\operatorname{Pic}_{X_n/A_n}^\circ$ simple sur $A_n$, et se borner à prouver
que $\operatorname{Pic}_{X_{n+1}/A_{n+1}}^\circ$ est simple sur $A_{n+1}$. On
constate qu'il suffit pour ceci de construire un schéma abélien $P_{n+1}$ sur
$A_{n+1}$ qui prolonge
\[
P_n=\operatorname{Pic}_{X_n/A_n}^\circ,
\]
et en même temps un Module inversible $\mathcal L_{n+1}$ sur
$X_{n+1}\times_{A_{n+1}}P_{n+1}$ qui prolonge le Module inversible
$\mathcal L_n$ sur $X_n\times_{A_n}P_n$ qui intervient dans la définition du
schéma de Picard $\operatorname{Pic}_{X_n/A_n}$ comme solution d'un problème
universel. (N.B. --- Nous pouvons supposer que $X$ est muni d'une section sur
$S$~...) Pour faire la construction, on doit utiliser le résultat-clef suivant :
tout schéma abélien défini sur un quotient d'anneau local artinien peut se
prolonger (en d'autres termes, le « schéma formel des Modules » absolu ([3],
II) pour un schéma abélien sur un corps algébriquement clos est simple sur
l'anneau des

\src{240}
vecteurs de Witt sur $k$) ; ce résultat lui-même s'obtient en utilisant les
propriétés formelles générales de l'obstruction à relever, et les opérations du
groupe. Ce résultat admis, on commence par prolonger $P_n$ n'importe comment en
$P_{n+1}$ ; on rencontre alors une obstruction à relever $\mathcal L_n$,
obstruction qui se trouve dans
\[
H^2(X_0\times P_0,\mathcal O_{X_0\times P_0})\otimes_k V
\qquad(\text{où }V=\mathfrak m^{n+1}/\mathfrak m^{n+2}),
\]
et même dans le sous-espace
\[
H^1(X_0,\mathcal O_{X_0})\otimes
H^1(P_0,\mathcal O_{P_0})\otimes_k V
\]
(compte tenu que la restriction de $\mathcal L_n$ aux deux facteurs $X_n$ et
$P_n$ est triviale). Or ce dernier espace n'est autre que
\[
H^1(P_0,\mathcal T_{P_0/k})\otimes V,
\]
où $\mathcal T_{P_0/k}$ est le faisceau tangent à $P_0/k$, donc c'est aussi
l'espace qui exprime l'indétermination qu'il y avait dans le relèvement de
$P_n$ à $P_{n+1}$ ([3], II). Par suite, on peut corriger ce relèvement (d'une
et d'une seule façon comme il se doit) de façon à tuer l'obstruction à relever
$\mathcal L_n$.

\result{COROLLAIRE}{3.6} Sous les conditions de (3.5),
$R^1f_*(\mathcal O_X)$ est un Module localement libre sur $S$ au voisinage de
$s$, et sa formation commute au changement de base.

Ce Module n'est autre en effet que le Module tangent à
$\operatorname{Pic}_{X/S}$ le long de la section unité.

\result{REMARQUE}{3.7} Utilisant le même argument que pour (3.5), on peut
montrer que si $S'$ est un sous-schéma de $S$ défini par un Idéal cohérent
nilpotent, et si on suppose seulement que
$\operatorname{Pic}_{X'/S'}^\circ$ existe et est à fibres simples, alors
$\operatorname{Pic}_{X/S}^\circ$ existe nécessairement et est un schéma
abélien sur $S$. Cela permet de construire des schémas de Picard dans certains
cas malgré l'absence d'hypothèse projective ; par exemple le schéma abélien
dual d'un schéma abélien sur un anneau artinien existe toujours. D'autre part,
utilisant (3.6) dans le cas où $X$ est un schéma abélien sur $S$, et utilisant
la structure connue de $H^*(X_s,\mathcal O_{X_s})$ comme algèbre extérieure
sur $H^1(X_s,\mathcal O_{X_s})$ (Rosenlicht--Serre), on trouve que pour tout
$i$, $R^if_*(\mathcal O_X)$ est localement libre, et de façon précise isomorphe
à la puissance extérieure $i$-ième de $R^1f_*(\mathcal O_X)$.

\result{REMARQUE}{3.8} Dans le cas d'un morphisme projectif simple
$f:X\longrightarrow S$, $S$ réduit et à caractéristiques résiduelles nulles,
le résultat (3.6) était connu par voie transcendante, comme conséquence de la
théorie de Hodge. En fait, tous les $R^pf_*(\Omega^q_{X/S})$ sont alors
localement libres. On a des contre-exemples par contre

\src{241}
dans le cas d'inégales caractéristiques pour « $R^1f_*(\mathcal O_X)$
localement libre », à l'aide des variétés de Serre ([8], p.~50). Il semble
qu'on n'ait aucun contre-exemple en égale caractéristique.

\result{COROLLAIRE}{3.8} Sous les conditions de (3.5),
$\operatorname{Pic}_{X/S}|U$ est simple sur $U$ en tous les points de
$\operatorname{Pic}_{X/S}^\sigma|U$.

On applique (2.10), (i). En particulier, tenant compte de (2.10), (iii) :

\result{COROLLAIRE}{3.9} Supposons $S$ à caractéristiques résiduelles nulles.
Alors $\operatorname{Pic}_{X/S}^\tau$ est simple sur $S$.

\result{REMARQUE}{3.10} On en déduit par exemple, si $\operatorname{Pic}_{X/S}^\tau$ est propre sur $S$, que le groupe de torsion
de Néron--Severi des fibres géométriques de $f$ reste constant sur toute
composante connexe de $S$ (ce qui est d'ailleurs évident par voie transcendante
lorsque $f$ est simple projectif). Signalons que l'utilisation directe de (2.5) permet de montrer, plus généralement, que, si $\operatorname{Pic}_{X/S}^\tau$ est propre sur $S$ (par exemple si $f:X\longrightarrow S$ est simple et projectif) et si $q$ est un nombre premier
distinct des caractéristiques résiduelles de $S$, alors la composante
$q$-primaire des groupes de torsion de Néron--Severi des fibres géométriques de
$X/S$ reste constante sur toute composante connexe de $S$. Il n'en est
cependant plus de même en caractéristique $p>0$ pour la composante $p$-primaire
de la torsion. Il reste cependant possible que le rang sur $k(s)$ de
\[
\operatorname{Pic}_{X_s/k(s)}^\tau/
\operatorname{Pic}_{X_s/k(s)}^{\circ\circ}=T_{X_s/k(s)}
\]
reste localement constant ; lorsque $S$ est réduit, on peut montrer qu'il
revient au même que $\operatorname{Pic}_{X/S}^{\circ\circ}$ existe, et que
$\operatorname{Pic}_{X/S}^\tau$ est plat sur $S$, et il suffit de tester la
question pour $S$ spectre d'un anneau de valuation discrète. C'est ce que j'ai
vérifié dans les quelques exemples que j'ai regardés ; mais comme l'énoncé
correspondant avec $S$ artinien est faux (cf. remarque (2.9) et (3.4)), il ne
faut sans doute pas se faire des illusions excessives.

\fsection{4}{Le théorème de finitude pour le schéma de Picard}

Soit $f:X\longrightarrow S$ un morphisme projectif et plat tel que
$\operatorname{Pic}_{X/S}$ existe, alors la considération des « polynômes de
Hilbert » $Q\in\mathbf Q[t]$ permet de décomposer
$\operatorname{Pic}_{X/S}$ en somme d'ouverts
$\operatorname{Pic}_{X/S}^Q$. Lorsqu'on ne fait pas des hypothèses
supplémentaires, assurant par exemple que $\operatorname{Pic}_{X/S}$ est séparé
sur $S$, il ne sera pas vrai en général que ces ouverts soient de type fini sur
$S$ ; on obtient un contre-exemple lorsque $X$ est « une conique dégénérant en
deux droites ». Il est

\src{242}
possible qu'il en soit cependant ainsi si $f$ est séparable et à fibres
géométriques irréductibles. La question est liée à celle de savoir si
$\operatorname{Pic}_{X/S}^\tau$ est de type fini sur $S$, ce qui est peut-être
vrai sans aucune hypothèse sur les fibres de $X/S$. Lorsque
$f:X\longrightarrow S$ est simple, on note que
$\operatorname{Pic}_{X/S}^\tau$ est contenu dans un des
$\operatorname{Pic}_{X/S}^Q$ (grâce au fait que, sur une variété projective
non singulière, l'équivalence « de torsion » est plus fine que (en fait
« identique à », d'après Matsusaka) l'équivalence numérique pour les
diviseurs), donc est de type fini sur $S$ si les
$\operatorname{Pic}_{X/S}^Q$ le sont. Noter que les questions de finitude qu'on
vient de soulever gardent un sens évident sans supposer même l'existence de
$\operatorname{Pic}_{X/S}$, car elles s'expriment en disant que certaines
familles de Modules inversibles sont « limitées » au sens de [3], IV :
considérons, pour tout corps algébriquement clos $k$, les sous-schémas intègres
de $\mathbf P_k^r$, de dimension $n$ et de degré $d$, et les Modules
inversibles sur de tels préschémas, ayant un polynôme de Hilbert $Q$ (où
$r,n,d,Q$ sont donnés) ; montrer que la famille de ces Modules (considérés
comme des Modules cohérents sur les $\mathbf P_k^r$) est limitée, i.e. peut se
paramétrer par un schéma de type fini sur $\mathbf Z$.

Utilisant la méthode de Matsusaka [6], une démonstration assez technique
(faisant appel aux « critères d'équivalence » sous une forme convenable) permet
de répondre par l'affirmative lorsque l'on se borne aux sous-variétés non
singulières de $\mathbf P_k^r$. De façon plus précise, on obtient le résultat
suivant :

\result{THÉORÈME}{4.1} Soit $f:X\longrightarrow S$\footnote{Le texte imprimé
porte « Soient : $X\longrightarrow S$ » ; le morphisme utilisé dans la suite
est $f$.} un morphisme simple projectif à fibres géométriques connexes, avec
$S$ noethérien, $\mathcal O_X(1)$ un Module très ample sur $X$ relatif à $S$,
$E$ une partie de $\operatorname{Pic}_{X/S}$, correspondant à une famille
$(\mathcal L_i)$ de Modules inversibles sur les fibres géométriques
$\overline X_{s_i}$ de $X/S$, $D_i$ un diviseur (pas nécessairement positif)
sur $\overline X_{s_i}$ définissant $\mathcal L_i$,
\[
a_n^{(i)}t^n+\cdots+a_0^{(i)}
\]
le polynôme de Hilbert de $\mathcal L_i$, $\xi=\xi_i$ un diviseur définissant
$\mathcal O_X(1)$, i.e. une section hyperplane. Les conditions sont
équivalentes :

\noindent a. $E$\footnote{Le texte imprimé porte ici « $Q$ » ; la partie
introduite dans l'énoncé et dont la quasi-compacité exprime que la famille est
limitée est $E$.} est quasi-compact, i.e. contenu dans un ouvert de type fini
sur $S$, i.e. la famille des $(\mathcal L_i)$ est limitée.

\noindent b. $E$ est contenu dans la réunion d'un nombre fini des ensembles
$\operatorname{Pic}_{X/S}^Q$, $Q\in\mathbf Q[t]$, i.e. l'ensemble des
polynômes de Hilbert des $\mathcal L_i$ est fini.

\noindent b'. (Si les fibres de $X/S$ ont toutes même dimension $n$.) Les deux
coefficients $a_{n-1}^{(i)}$ et $a_{n-2}^{(i)}$ des polynômes de Hilbert des
$\mathcal L_i$ restent dans un ensemble fini.

\src{243}

\result{COROLLAIRE}{4.2} Soit $f:X\longrightarrow S$ un morphisme simple
projectif à fibres géométriques connexes, alors les schémas
$\operatorname{Pic}_{X/S}^Q$ et $\operatorname{Pic}_{X/S}^\tau$
($Q\in\mathbf Q[t]$) sont projectifs sur $S$.

Comme ils sont de type fini sur $S$ d'après (4.1), on peut appliquer (2.4).

\section*{BIBLIOGRAPHIE}
\addcontentsline{toc}{section}{Bibliographie}

\begin{description}
\item[{[1]}] CHEVALLEY (Claude). -- Sur la théorie de Picard,
\emph{Amer. J. of Math.}, t. 82, 1960, p. 435--490.

\item[{[2]}] GROTHENDIECK (A.) et DIEUDONNÉ (J.). -- Éléments de géométrie
algébrique, Chapitre I et suivants. -- Paris, Presses universitaires de France,
1960--1961 (Institut des Hautes Études Scientifiques, 4, 8, 11, \ldots).

\item[{[3]}] GROTHENDIECK (Alexander). -- Technique de descente et théorèmes
d'existence en géométrie algébrique, I--V, \emph{Séminaire Bourbaki}, t. 12,
1959/60, n\textsuperscript{o} 190, 29 p. et n\textsuperscript{o} 195, 22 p. ;
t. 13, 1960/61, n\textsuperscript{o} 212, 20 p. et
n\textsuperscript{o} 221, 28 p. ; t. 14, 1961/62,
n\textsuperscript{o} 232, 19 p.

\item[{[4]}] GROTHENDIECK (Alexander). -- Séminaire de géométrie algébrique,
1960/61. -- Paris, Institut des Hautes Études Scientifiques, 1961.

\item[{[5]}] GROTHENDIECK (Alexander). -- Séminaire de géométrie algébrique,
1961/62, rédigé par un groupe d'auditeurs. -- Paris, Institut des Hautes Études
Scientifiques (à paraître).

\item[{[6]}] MATSUSAKA (T.). -- The criteria for algebraic equivalence and the
torsion group, \emph{Amer. J. Math.}, t. 79, 1957, p. 53--66.

\item[{[7]}] MUMFORD (D.) and TATE (J.). -- Séminaire de géométrie algébrique,
Harvard University, Spring Term 1962 (à paraître).

\item[{[8]}] SERRE (Jean-Pierre). -- Sur la topologie des variétés algébriques
en caractéristique $p$, \emph{Symposium internacional de topología algebraica}
[1956, Mexico], p. 24--53. -- Mexico, Universidad nacional autónoma, 1958.
\end{description}

\par\smallskip
\begingroup\footnotesize
\noindent\textit{La clause b'' imprimée reste consignée dans C086; l'exclut de la liste d'équivalences du texte courant à cause du contre-exemple Stacks D000571.}\par
\noindent\textit{à 0058. En 3.2, la lecture imprimée suppose $X$ normal et propre sur $k$, sans normalité géométrique (C122). En 3.10, les deux hypothèses imprimées de propreté portent sur le schéma de Picard entier, et non sur $\operatorname{Pic}_{X/S}^{\tau}$ (C124). C123 rétablit le $\sigma$ imprimé en 2.10(i) et 3.8 : le cercle et $\tau$ étaient des erreurs de transcription, non des lectures imprimées.}\par
\endgroup
